\clearpage
\fancyhead[L]{\small CONNAISSANCE DES TEMS.}
\hypertarget{native-497-lower}{}\label{native-497-lower}
\begin{center}
\small Göttingische gelehrte Anzeigen. Stück 56. 57. S. 553..562. 7. April 1808.\\[3pt]
\rule{.72\linewidth}{.3pt}
\end{center}

\emph{Connoissance des tems, ou des mouvemens célestes à l'usage des astronomes et des navigateurs pour l'an} 1808, publiée par le bureau des longitudes. Paris. De l'imprimerie impériale. 1806. 502 Seiten im Oktavformat. --- Dieselbe für das Jahr 1809. Ebenda. 1807. 502 Seiten im Oktavformat.

Mit diesen Jahrgängen der zu Recht so hoch geschätzten astronomischen Ephemeriden ist man zur gewöhnlichen Zeitrechnung zurückgekehrt, nachdem 12 Jahrgänge (vom Jahr IV bis XII) in der Form des nun vergessenen republikanischen Kalenders erschienen waren. Die auch hier unverändert beibehaltene Anordnung des astronomischen Kalenders selbst und der damit verbundenen regelmäßigen Beiträge ist zu bekannt, als dass wir uns dabei aufhalten müssten. Nur bei der Tafel zur geografischen Lage der wichtigsten Orte der Erde --- die gewöhnlich von Zeit zu Zeit einige Ergänzungen und Verbesserungen erhält, obwohl selbst ganz bekannte Bestimmun%
\SourceLeaf{498}%
gen noch nicht überall berücksichtigt sind und manche groben Druckfehler von einem Jahrgang zum nächsten fortgeschleppt werden --- bemerken wir die mit dem Jahr 1809 eingeführte Änderung. Die Tafel ist nicht mehr wie bisher nach Ländern geordnet, sondern läuft für die ganze Erde in einer einzigen alphabetischen Folge durch. Der Rezensent würde der bisherigen Anordnung den Vorzug vor dieser Neuerung geben, da es mitunter bequem ist, die vorliegenden guten Ortsbestimmungen für ein Land beisammen zu haben. Wir wenden uns also sogleich den \emph{Zusätzen} zu, die stets aus einer großen Zahl wertvoller Abhandlungen, Beobachtungen und interessanter Mitteilungen bestehen und den Ephemeriden für Astronomen bleibenden Wert verleihen.

Die Redaktion der Zusätze, die bisher \textsc{Lalande} besorgt hatte, hat vom Jahr 1808 an \textsc{Delambre} übernommen. Zugleich hat man in diesem Jahr damit begonnen, sie mit den Beobachtungen von \textsc{Bouvard} an der Pariser kaiserlichen Sternwarte zu eröffnen, und wird dies künftig beibehalten. Diese regelmäßige Veröffentlichung der ebenso zahlreichen wie wertvollen Beobachtungen, auf dieselbe Weise, wie die Beobachtungen aus Greenwich bekannt gemacht werden, muss den Astronomen sehr willkommen sein und erhöht den Wert der Connoissance des tems beträchtlich. Im Jahrgang 1808 nehmen die Beobachtungen vom 24. September 1803 bis zum 21. September 1804 nach einer kurzen Beschreibung der Instrumente 83 Seiten ein. Das vortreffliche Passageninstrument von \textsc{Berge} hat $2\frac{1}{2}$ Meter (92 Pariser Zoll) Brennweite und 11 Zentimeter (4 Zoll) Öffnung und wurde im September 1803 aufgestellt. Die Uhr ist von \textsc{Louis Berthoud}. Der acht Fuß lange Mauerquadrant von \textsc{Bird}, der mit einem achromatischen Fernrohr ausgestattet ist, ist derselbe, der früher \textsc{Lemonnier} gehörte. Unter den Beobachtungen finden wir auch zahlreiche bisher noch nicht verwendete Beobachtungen der Ceres und Pallas. Angefügt sind Beobachtungen und parabolische Elemente des Kometen von 1804 und der beiden von 1805, von letzteren jedoch nur die reduzierten Längen und Breiten. An dieses astronomische Journal von \textsc{Bouvard} schließen sich Beobachtungen von \textsc{Flaugerques} in Viviers und von \textsc{Vidal} in Toulouse an. Letztere bestehen aus Kulminationen von Merkur und Venus, die dieser scharfsichtige Astronom zum Teil in derselben Minute wie die Sonne beobachtet hat. Sonnenfinsternisse sowie Planeten- und Sternbedeckungen, beobachtet von \textsc{Scarpellizi} in Rom und zum Teil von \textsc{Lalande} berechnet. Von \textsc{Lalande} Beobachtungen der unteren Konjunktionen der Venus von 1802 und 1804 und der Sonnenfinsternis vom 1. Februar 1804 von mehreren Orten aus. Von \emph{ebendemselben} ein neues Verzeichnis der Eigenbewegung von 500 Sternen. Der Rezensent zählt darunter 28, bei denen die jährliche Eigenbewegung in Rektaszension, und 15, bei denen sie in Deklination eine halbe Bogensekunde übersteigt. Es scheint aber, dass \textsc{Lalande} die Ungleichförmigkeit der Präzession nicht gebührend berücksichtigt hat, obwohl er dies ausdrücklich versichert. Wenigstens finden wir, dass dies beim Polarstern nicht geschehen ist. Vielmehr ist aus der Position von \textsc{Lacaille} von 1750 und der von \textsc{von Zach} von 1800 unmittelbar auf die beobachtete jährliche Bewegung in Rektaszension $173''18$ geschlossen worden. Diese ist mit der für 1800 (nicht, wie es sein sollte, für 1775) berechneten Bewegung $207''52$ verglichen worden. Danach wird die jährliche Eigenbewegung ganz falsch mit $34''34$ angesetzt. --- Von \textsc{Prony} eine Abhandlung über die Berechnung der geografischen Längen und Breiten aus den Abständen vom Meridian und von der Senkrechten dazu auf dem elliptischen Sphäroid. Seine Formeln sind zwar nur auf mäßige Abstände anwendbar, empfehlen sich aber durch ihre Einfachheit. Von \textsc{Delambre} eine neue, einfallsreiche Methode, die Anordnung der Jupitermonde allein durch Rechnung mithilfe bestimmter Hilfstafeln zu finden, wozu bisher gewöhnlich mechanische Hilfsmittel verwendet wurden. Von \emph{demselben} die Geschichte der Astronomie für 1804 und 1805. Die Entdeckung der \emph{Juno} nimmt hier zu Recht den ersten Platz ein. Das Urteil, dass sich jeder Astronom stets glücklich schätzen würde, eine so große Entdeckung mit einem ganzen Jahr --- vor der Entdeckung der Ceres hätte sicher niemand gezögert zu sagen, mit einem ganzen Leben --- angestrengter Arbeit zu erkaufen, macht dem französi%
\SourceLeaf{499}%
schen Astronomen Ehre. Dass die Astronomie in sechs Jahren viermal das Glück hatte, eine so unschätzbare Bereicherung zu erfahren, kann deren Wert wahrlich nicht mindern. Vielmehr müssen diese vor allem für Deutschland so ruhmvollen Ereignisse den Mut und die Tatkraft der Astronomen umso mehr anspornen, da sie die Hoffnung auf ähnliche künftige Erfolge so sehr erhöhen. Umso befremdlicher und mit jenem Urteil ganz unvereinbar ist die Kälte, mit der sich \textsc{Delambre} über den Plan äußert, den vor sieben Jahren, noch vor der Entdeckung der Ceres, in Lilienthal vor allem Herr \textsc{von Zach} entworfen hatte. Der ganze Himmel sollte unter einer Anzahl von Astronomen aufgeteilt und fortlaufend abgesucht werden. Auch er war zur Teilnahme eingeladen. Er hält einen solchen Vergleich für höchst beschwerlich und langweilig, weil man jede Nacht stundenlang mechanisch am unbeweglichen Fernrohr stehen und Sterne zählen müsste. Er glaubt, dass dies jeden ganz abschrecken müsse, der sich zur Astronomie berufen fühle. Nur etwa \emph{beiläufig}, sooft man den Mond, die anderen Planeten und die größeren Fixsterne beobachte, solle man einen glücklichen Zufall abwarten. Nur der \emph{verdiene} eine solche Entdeckung zu machen, dem sie gelinge, ohne darauf ausgegangen zu sein, während er mit einem großen und \emph{wirklich nützlichen} Vorhaben beschäftigt sei, und so weiter. Äußerungen dieser Art von einem Mann, der wegen seiner vielen verdienstvollen Arbeiten so große Achtung verdient, darf man nicht mit Stillschweigen übergehen. Männer. deren Name als Autorität gilt, sollten sich, meinen wir, umso mehr vor Urteilen hüten, die im Geist unseres engherzigen, frivolen Jahrhunderts gesprochen zu sein scheinen könnten. Man glaubt in diesem Jahrhundert die Musen zu ehren, indem man sie zu bloßen Mägden unserer Bedürfnisse herabwürdigt. Oder man schätzt, wie einer unserer geistreichsten Schriftsteller sich ausdrückt, in der erhabensten Astronomie die Gestirne nur als Schrittzähler und Wegweiser für Pfefferflotten. Wer nicht genug Begeisterung für die Wissenschaft empfindet, um das Große an der Entdeckung eines bisher unbekannten Himmelskörpers zu fühlen, der wie die Erde regelmäßig unsere Sonne umkreist und mit dessen Erscheinungen sich Himmelsbeobachter noch nach Jahrtausenden beschäftigen werden, und wessen erster Gedanke dabei ist, wozu eine solche Entdeckung nützen werde, dem haben wir nichts weiter zu sagen als unsere Überzeugung, dass gerade die neuen Planeten der theoretischen und physischen Astronomie einen ganz neuen Aufschwung geben werden. Wer aber eine andere Ansicht hat und in den vom glücklichsten Erfolg gekrönten Nachforschungen unserer \textsc{Olbers} und \textsc{Harding} einen Grund zu der keineswegs wirklichkeitsfremden Hoffnung sieht, dass noch manche ähnlichen Erfolge zu erwarten sind, der muss lebhaft wünschen, dass jener Plan bald ausgeführt wird. Er allein kann uns in einem für einzelne Astronomen nicht zu umfangreichen Unternehmen vom Zufall unabhängig machen. Ein solcher Leser wird sehr gern hören, dass ein wichtiger Schritt dazu schon getan ist. Mit wahrer Freude berichten wir, dass die Himmelskarten unseres Herrn Professor \textsc{Harding}, die das Unternehmen so vortrefflich vorbereitet haben, trotz der für die Wissenschaften bei uns so bedrängten Zeitumstände kurz vor dem Erscheinen stehen. Es ist übrigens eine ganz unbegründete Vorstellung von \textsc{Delambre}, dass die Arbeit bei solchen Nachforschungen so beschwerlich, zeitraubend und langweilig sei. Werden sie nur auf zweckmäßige Weise angestellt, nicht mit feststehenden Instrumenten, sondern mit guten lichtstarken Nachtfernrohren, so kann man in kurzer Zeit ein beträchtliches Gebiet absuchen, sobald einmal vollständige Karten vorliegen. In dieser Beschäftigung wird man gewiss immer eine angenehme Erholung von anderen mechanischen Arbeiten der beobachtenden und rechnenden Astronomie finden. Auch ist es keineswegs nötig, diese Durchmusterung jede Nacht vollständig zu wiederholen.

Nach dieser Abschweifung, die uns als ein Wort zur rechten Zeit erschien, kommen wir auf die übrigen Beiträge jenes Aufsatzes zurück. Über die Kometen von 1804 und 1805. \textsc{Schröter}s und \textsc{Harding}s
\SourceLeaf{500} Beobachtungen des Saturnrings. \textsc{Schröter}s und \textsc{Herschel}s Bestimmungen der Durchmesser der neuen Planeten. \textsc{Piazzi}s Beobachtungen der Parallaxe einiger Fixsterne. Einige den Berliner astronomischen Jahrbüchern entnommene Formeln von \textsc{Camerer} und \textsc{Olbers}. Über die Verfahren von \textsc{Bowditch} und \textsc{Mendoza}, scheinbare Mondabstände auf die wahren zu reduzieren. Mitteilungen aus dem zweiten Band der \emph{Asiatic Researches} über die Methoden der Hindus zur Berechnung der Sinustafeln. Die übrigen kürzeren Beiträge dieses Jahrgangs enthalten Anzeigen verschiedener neuer astronomischer Werke, die amtlichen Verhandlungen über die Abschaffung des republikanischen und die Wiedereinführung des gregorianischen Kalenders, Beobachtungen der Sonnenfinsternis vom 16. Juni 1806 und des letzten Merkurdurchgangs, die Ausschreibung des Preises für die Berechnung der Störungen der Pallas, einige Verbesserungen an den Tafeln in \textsc{Lalande}s Astronomie und einen Auszug aus den meteorologischen Beobachtungen in Paris aus dem Jahr XII.

Im \emph{Jahrgang} für 1809 nehmen die an der kaiserlichen Sternwarte vom 23. September 1804 bis zum Ende des Jahres 1805 angestellten Beobachtungen 108 Seiten ein. Die neuen Planeten Ceres, Pallas und Juno wurden 1804 fleißig im Meridian beobachtet, die letzte bis zum 24. Dezember. Doch bedauern wir, dass Beobachtungen von Pallas und Ceres vom Ende des Jahres 1805 fehlen. Von den beiden Kometen des Jahres 1805 finden wir hier die Originalbeobachtungen. Es folgen die astronomischen Beobachtungen \textsc{Messier}s aus fünf Jahren, von 1760\mbox{. . . .}1764, die sich an die bereits in früheren Bänden der Connoissance des tems mitgeteilten anschließen. Man findet hier eine Sonnenfinsternis, zwei Mondfinsternisse, mehrere Bedeckungen von Fixsternen, auch eine des Mars und sehr viele Verfinsterungen der Jupitermonde. Für Freunde der Meteorologie heben wir die Bemerkung hervor, dass am 19. Juli 1760 zwischen vier und fünf Uhr nachmittags das Quecksilberthermometer in Paris auf 32 Grad stieg. --- Astronomische Beobachtungen in Viviers von \textsc{Flaugerques}. --- Von \textsc{Burckhardt} die Anordnung des Räderwerks zur Darstellung der Planetenbewegungen, einige Bemerkungen dazu, wie man den Gebrauch des \textsc{Borda}schen Kreises bei verschiedenen Gelegenheiten genauer, allgemeiner und bequemer machen kann, sowie Untersuchungen zur Eigenbewegung unseres Sonnensystems anhand eines Vergleichs der beobachteten Eigenbewegungen der Fixsterne. Das Ergebnis ist, dass sich darüber noch wenig Befriedigendes feststellen lässt. Einige ältere chinesische Beobachtungen von Sonnenwenden, aus einem Manuskript, das Pater \textsc{Gaubil} einst an \textsc{Delisle} geschickt hat und das in der kaiserlichen Bibliothek aufbewahrt wird. Ebenfalls von \textsc{Burckhardt} der Vorschlag einer neuen Bauweise für Spiegelteleskope. Er fängt die vom großen Spiegel zurückgeworfenen Strahlen auf halber Brennweite mit einem senkrecht zur Achse stehenden Planspiegel auf, der sie zu den Okularen im durchbohrten Hohlspiegel zurückschickt. Diese Bauweise hat gegenüber der von \textsc{Newton} den Vorteil, dass das ganze Teleskop bei gleicher Brennweite des Spiegels nur halb so lang wird. Dafür entzieht der kleine Spiegel, dessen Durchmesser halb so groß sein muss, dem großen ein Viertel des in den Tubus einfallenden Lichts. Der Instrumentenbauer \textsc{Caroché} wird versuchen, den Vorschlag auszuführen. --- Bemerkungen über eine von \textsc{Ducum} in Bordeaux vorgeschlagene Methode, die Breite auf See durch zwei außerhalb des Meridians gemessene Höhen zu bestimmen. Der Berichterstatter hat diesen Bemerkungen noch verschiedene eigene Zusätze zu diesem Problem und der Lösung von \textsc{Douwes} beigefügt. --- Von \textsc{Oltmanns} eine wertvolle Untersuchung über die Länge von Quito anhand der von Herrn \textsc{v. Humboldt} beobachteten Mondabstände und Verfinsterungen der Jupitermonde sowie verschiedener älterer Beobachtungen. Sein Ergebnis für diese bisher so unterschiedlich bestimmte Länge ist $5^{\mathrm u}\,24^{\mathrm m}\,20^{\mathrm s}$ westlich von Paris. --- Über eine Stelle in der Geografie des Ptolemäus zur Projektionsart des Marinus von Tyrus, in der man ganz zu Unrecht schon eine Spur von Mercators Projektion zu erkennen glaubte. --- Beobachtungen von Ceres, Pallas und Juno durch \textsc{Poczobut} und \textsc{Keschka} in Wilna. --- Von \textsc{Henry} strenge Formeln für die Parallaxe der Länge und
\SourceLeaf{501} Breite. --- Eine neue Art Mikrometer von \textsc{Prony}. Die zugrunde liegende Idee ist ebenso einfach wie einfallsreich und lässt sich auch bei anderen Gelegenheiten mit Nutzen anwenden. Um nämlich dem beweglichen Faden (dem Läufer) eine sehr langsame Bewegung zu geben, ohne dazu eine Schraube mit sehr feinem Gewinde zu benötigen, die sich schwer herstellen lässt und sich bald verändert, bewegt er den Faden durch eine Schraube mit ziemlich grobem Gewinde. Auf beiden Seiten dieses Gewindes gibt er der Schraube aber auf derselben Achse weitere Gewindegänge, die ein wenig enger oder weiter als die mittleren und untereinander gleich sind. Die zugehörigen Muttern sitzen fest in der Büchse. Auf diese Weise hängt die tatsächliche Bewegung des Läufers nur vom Unterschied der beiden Arten von Gewindegängen ab. --- Auszüge aus einigen neuen astronomischen Schriften. Interessant sind die Auszüge aus einer Abhandlung von \textsc{Monteiro} über die Berechnung der Finsternisse im vierten Band der Ephemeriden von Coimbra. --- Geschichte der Astronomie für 1807. Sie betrifft hauptsächlich die Entdeckung der Vesta. Deutsche Leser werden darüber hier nichts Neues finden. Vom neuesten Kometen finden wir hier noch die ersten Beobachtungen von \textsc{Thulis} vom 21. und 22. September. Die Deklinationen sind nur in ganzen Minuten angegeben. --- Einige geografische Bestimmungen im Mittelmeer und an dessen Küste von \textsc{Dionisio Alcala Galeano}. --- Diesen Jahrgang beschließt wieder ein Auszug aus \textsc{Bouvard}s meteorologischen Beobachtungen aus dem Jahr XIII.
